\documentclass[12pt]{article}
\usepackage{amsmath,amssymb,amsthm}
\usepackage{graphicx}
\usepackage{tikz}
\usetikzlibrary{calc}
\usepackage{booktabs}
\usepackage[shortlabels]{enumitem}
\usepackage[colorlinks=true,linkcolor=blue!45!black,citecolor=blue!45!black,urlcolor=blue!45!black]{hyperref}
\hypersetup{pdftitle={Maximum Blackouts That Keep Every Lattice Rectangle Identifiable},pdfauthor={Arjun Pemmasani}}

\definecolor{gridcol}{gray}{0.86}
\definecolor{keptcol}{RGB}{226,236,247}
\definecolor{keptline}{gray}{0.35}
\definecolor{hidcol}{gray}{0.62}
\definecolor{pencol}{gray}{0.55}
\definecolor{cyccol}{RGB}{31,94,168}
\definecolor{rectA}{RGB}{0,114,178}
\definecolor{rectB}{RGB}{213,94,0}
\definecolor{rectC}{RGB}{0,158,115}
\definecolor{excesscol}{RGB}{253,214,160}
\tikzset{
  lab/.style={font=\scriptsize,text=black!55},
  capt/.style={font=\small},
  tiny/.style={font=\scriptsize},
  rowv/.style={draw,rectangle,fill=white,inner sep=1.5pt,minimum size=13pt,font=\scriptsize},
  colv/.style={draw,circle,fill=white,inner sep=0.5pt,minimum size=14pt,font=\scriptsize},
}

\theoremstyle{plain}
\newtheorem{theorem}{Theorem}
\newtheorem{lemma}[theorem]{Lemma}
\newtheorem{proposition}[theorem]{Proposition}
\newtheorem{corollary}[theorem]{Corollary}
\newtheorem{conjecture}[theorem]{Conjecture}
\newtheorem*{thmA}{Theorem A}
\newtheorem*{thmB}{Theorem B}
\newtheorem*{thmC}{Theorem C}
\newtheorem*{thmD}{Theorem D}
\newtheorem*{thmE}{Theorem E}
\newtheorem*{thmF}{Theorem F}
\theoremstyle{definition}
\newtheorem{definition}[theorem]{Definition}
\newtheorem{example}[theorem]{Example}
\theoremstyle{remark}
\newtheorem{remark}[theorem]{Remark}

\newcommand{\seqnum}[1]{\href{https://oeis.org/#1}{\rm \underline{#1}}}
\newcommand{\pres}{P_S}
\newcommand{\kA}{\mathsf{A}}
\newcommand{\kB}{\mathsf{B}}
\newcommand{\kC}{\mathsf{C}}
\newcommand{\kE}{\mathsf{E}}

\title{Maximum Blackouts That Keep Every Lattice Rectangle Identifiable}
\author{Arjun Pemmasani\\
Department of Mathematics\\
Harvey Mudd College\\
Claremont, CA 91711\\
USA\\
\texttt{apemmasani@hmc.edu}}
\date{}

\begin{document}
\maketitle

\begin{abstract}
Kagey's Problem 001 asks how many points of a rectangular grid of points can be
blacked out so that every rectangle with its four corners on grid points, tilted
rectangles included, can still be told apart from every other rectangle by its white
corners alone. We study the largest size of such a blackout and the number of
blackouts of that size, the tables A397315 and A399580 in the On-Line Encyclopedia
of Integer Sequences. We prove the exact maximum on two, three and four rows, count
the three-row optima exactly, and show that their number grows like a constant times
$m^2 3^m$ with an explicit rational generating function. On four rows
the optima are an exponentially small fraction of the candidates. A finite collar
theorem decides when a construction extends to every width, and it gives a five-row
family that refutes a conjecture recorded in the encyclopedia. For every fixed number
of rows the maximum minus the width is eventually periodic, and on squares the
maximum grows faster than linearly. The main tools are a hitting-set reduction, a
bipartite graph without four-cycles and six-cycles, and a classification of the tilted
rectangles that show at most two corners.
\end{abstract}

\input{sections/sec01_intro}
\input{sections/sec02_conv}
\input{sections/sec03_strip}
\input{sections/sec04_graph}
\input{sections/sec05_three}
\input{sections/sec06_four}
\input{sections/sec07_five}
\input{sections/sec08_height}
\input{sections/sec09_verify}
\input{sections/sec10_open}
\input{sections/sec11_ack}

\appendix
\input{sections/appA_fivecases}
\input{sections/appB_threegf}
\input{sections/appC_square}

\begin{thebibliography}{99}

\bibitem{alanko2011}
S. Alanko, S. Crevals, A. Isopoussu, P. R. J. \"Osterg\r{a}rd, and V. Pettersson,
Computing the domination number of grid graphs, \emph{Electron. J. Combin.}
\textbf{18} (2011), Paper P141.

\bibitem{alonspencer2016}
N. Alon and J. H. Spencer, \emph{The Probabilistic Method}, 3rd ed., Wiley, 2008.

\bibitem{babichev2026mfcs}
D. Babichev and S. Babichev, Counting all lattice rectangles in the square grid in
near-linear time, in \emph{51st International Symposium on Mathematical Foundations of
Computer Science (MFCS 2026)}, Leibniz International Proceedings in Informatics,
Vol.\ 386, Schloss Dagstuhl, 2026, pp.\ 26:1--26:17.

\bibitem{babichevpinchuk2026}
D. Babichev and D. Pinchuk, Computing all lattice-rectangle counts by rational
staircase sums, arxiv preprint arXiv:2607.17982 [cs.CG], 2026. Available at
\url{https://arxiv.org/abs/2607.17982}.

\bibitem{babichevshpakova2026}
D. Babichev and T. Shpakova, Counting lattice rectangles in $O(n\log n)$ operations,
arxiv preprint arXiv:2607.17961 [cs.CG], 2026. Available at
\url{https://arxiv.org/abs/2607.17961}.

\bibitem{baccelli1992}
F. Baccelli, G. Cohen, G. J. Olsder, and J.-P. Quadrat, \emph{Synchronization and
Linearity: An Algebra for Discrete Event Systems}, Wiley, 1992.

\bibitem{bellitto2018}
T. Bellitto, Separating codes and traffic monitoring, \emph{Theoret. Comput. Sci.}
\textbf{717} (2018), 73--85.

\bibitem{butler2019}
S. Butler, J. Choi, K. Kim, and K. Seo, Enumerating multiplex juggling patterns,
\emph{J. Integer Sequences} \textbf{22} (2019), Article 19.1.7.

\bibitem{conlonfoxsudakov2016}
D. Conlon, J. Fox, and B. Sudakov, Short proofs of some extremal results II,
\emph{J. Combin. Theory Ser. B} \textbf{121} (2016), 173--196.

\bibitem{decaenszekely1997}
D. de Caen and L. A. Sz\'ekely, On dense bipartite graphs of girth eight and upper
bounds for certain configurations in planar point-line systems, \emph{J. Combin.
Theory Ser. A} \textbf{77} (1997), 268--278.

\bibitem{erdos1080}
P. Erd\H{o}s, Problem 1080, in T. Bloom, ed., \emph{Erd\H{o}s Problems},
\url{https://www.erdosproblems.com/1080} (accessed October 2026).

\bibitem{euler2005}
R. Euler, The Fibonacci number of a grid graph and a new class of integer sequences,
\emph{J. Integer Sequences} \textbf{8} (2005), Article 05.2.6.

\bibitem{flajolet2009}
P. Flajolet and R. Sedgewick, \emph{Analytic Combinatorics}, Cambridge University
Press, 2009.

\bibitem{foucaud2015}
F. Foucaud, Decision and approximation complexity for identifying codes and
locating-dominating sets in restricted graph classes, \emph{J. Discrete Algorithms}
\textbf{31} (2015), 48--68.

\bibitem{glp2025}
J. Guilhot, E. Little, and J. Parkinson, Bounded weight functions on regular languages
and groups, arxiv preprint arXiv:2510.17061 [math.GR], 2025. Available at
\url{https://arxiv.org/abs/2510.17061}.

\bibitem{gyori1997}
E. Gy\H{o}ri, $C_6$-free bipartite graphs and product representation of squares,
\emph{Discrete Math.} \textbf{165/166} (1997), 371--375.

\bibitem{hardywright2008}
G. H. Hardy and E. M. Wright, \emph{An Introduction to the Theory of Numbers}, 6th ed.,
Oxford University Press, 2008.

\bibitem{jiang2018}
M. Jiang, Periodicity of identifying codes in strips, \emph{Inform. Process. Lett.}
\textbf{135} (2018), 77--84.

\bibitem{kageyG4G14}
P. Kagey, Open problem collection, Problem 1, Gathering 4 Gardner 14 gift exchange,
2020. Available at
\url{https://www.gathering4gardner.org/g4g14gift/G4G14-PeterKagey-OpenishProblems.pdf}.

\bibitem{kagey001}
P. Kagey, Problem 001, Open Problem Collection,
\url{https://peterkagey.com/problems/001/} (accessed October 2026).

\bibitem{katona1966}
G. Katona, On separating systems of a finite set, \emph{J. Combin. Theory}
\textbf{1} (1966), 174--194.

\bibitem{mantel1907}
W. Mantel, Problem 28 (solutions by H. Gouwentak, W. Mantel, J. Teixeira de Mattes,
F. Schuh, and W. A. Wythoff), \emph{Wiskundige Opgaven} \textbf{10} (1907), 60--61.

\bibitem{mertens2024}
S. Mertens, Domination polynomials of the grid, the cylinder, the torus, and the king
graph, \emph{J. Integer Sequences} \textbf{27} (2024), Article 24.7.6.

\bibitem{pulp}
S. Mitchell, M. O'Sullivan, and I. Dunning, PuLP: a linear programming toolkit for
Python, technical report, University of Auckland, 2011. Available at
\url{https://optimization-online.org/2011/09/3178/}.

\bibitem{naorverstraete2005}
A. Naor and J. Verstra\"ete, A note on bipartite graphs without $2k$-cycles,
\emph{Combin. Probab. Comput.} \textbf{14} (2005), 845--849.

\bibitem{neuwirth2001}
S. Neuwirth, The size of bipartite graphs with girth eight, arxiv preprint
arXiv:math/0102210 [math.CO], 2001. Available at
\url{https://arxiv.org/abs/math/0102210}.

\bibitem{oeis}
OEIS Foundation Inc., \emph{The On-Line Encyclopedia of Integer Sequences}, 2026.
Available at \url{https://oeis.org}.

\bibitem{trureturing2026}
The Omega Institute, Prove the exact Three-Row blackout maximum, pull request 8084 to
the \texttt{trureturing} repository, merged 15 September 2026. Available at
\url{https://github.com/the-omega-institute/trureturing/pull/8084}.

\bibitem{repo}
A. Pemmasani, Kagey-Problems repository, 2026, folder
\texttt{kagey-problems/problem001} of \url{https://github.com/apemm/Kagey-Problems}.

\bibitem{spencer1972}
J. Spencer, Tur\'an's theorem for $k$-graphs, \emph{Discrete Math.} \textbf{2} (1972),
183--186.

\end{thebibliography}

\noindent 2020 {\it Mathematics Subject Classification}: Primary 05A15; Secondary
05C35, 05B40, 52C05.

\noindent \emph{Keywords:} blackout, lattice rectangle, rectangle identification, test
cover, bipartite graph of girth eight, transfer matrix, rational generating function.

\noindent (Concerned with sequences
\seqnum{A001787},
\seqnum{A322284},
\seqnum{A397315}, and
\seqnum{A399580}.)

\end{document}
